\chapter{Structure projections}
\label{chap:proj}

\section*{At a glance}
\begin{itemize}
\item \lead{Goal} Represent \texttt{.proj S i e} as an application of $S$'s own recursor: no new
\texttt{VExpr} node, no new \texttt{IsDefEq} rule.
\item \lead{Main results} Theorem~\ref{thm:proj-instfields-minor-spine} (arithmetic core of the
iota reduct); Lemma~\ref{fam:proj-fn-subst} (projection stable under substitution);
Example~\ref{def:projshape-checkproj} (the real kernel accepts the expansion);
Example~\ref{thm:projinhabit-dep-hp1ty} (a dependent motive typed via one iota step, no eta).
\item \lead{Status} \stProved{} all 68 nodes. \stTainted{} Example~\ref{test:projinhabit-axioms}
pins an axiom list including \texttt{sorryAx}, from a \texttt{sorry} elsewhere
(\texttt{inferProj.WF\_struct}, Chapter~\ref{chap:trexpr}).
\item \lead{Contribution} Every node here is new. \texttt{trproj} adds the projection expansion,
its substitution metatheory, and its kernel- and model-side validation. \texttt{iota} adds the
executable tests deciding Chapter~\ref{chap:inductive}'s clauses against the kernel.
\item \lead{Lean files} \texttt{Lean4Lean/Theory/Proj.lean}, \texttt{Lean4Lean/Tests/ProjShape.lean},
\texttt{Lean4Lean/Tests/ProjInhabit.lean}, \texttt{Lean4Lean/Tests/IotaShape.lean},
\texttt{Lean4Lean/Tests/ShapeDecide.lean}.
\item \lead{Thesis} axioms.tex \S2.6; typesys.tex \S3.1 and Wtypes.tex \S5.1 for the projection
idiom (borrowed, not inherited: see Caveats).
\end{itemize}

\begin{itemize}
\item \lead{Where} contexts list outermost first, $\mathrm{bvar}\;0$ innermost: over
$\Gamma,f_0\dots f_{n-1}$ the field $f_j$ is $\mathrm{bvar}(n-1-j)$.
\item \lead{Where} the level list is per field ($uss\,j$): e.g.\ $\mathrm{Sigma.fst}$ needs
$\{u+1,u,v\}$ against $\mathrm{Sigma.snd}$'s $\{v+1,u,v\}$.
\end{itemize}

\section{The projection expansion}

Ten definitions build $P_i$ from a structure's field telescope, each layered on the last; one
lemma restates their unfolding as rewrite rules.

\begin{definition}[Field selector]
  \label{def:proj-field-selector}
  \lean{Lean4Lean.VExpr.fieldSelector}
  \leanok
  \uses{ind:vexpr}
  \contrib{trproj}\srcloc{Lean4Lean/Theory/Proj.lean}{70}
  \thesisref{axioms.tex \S2.6.3, the recursor: the minor premise
  $\varepsilon_c=\forall b{:}\beta.\forall v{:}\delta.\,C\,p[b]\,(c\,b)$.}

  \lead{In short} The minor premise that makes a structure's recursor \hl{return field $i$}.
  \begin{itemize}
  \item $\mathrm{fieldSelector}\;Fs\;i = Fs.\mathrm{foldr}\;\lambda\;(\mathrm{bvar}(|Fs|-1-i))$: a
        $\lambda$-telescope over the field types, returning the $i$-th binder from the outside.
  \item \lead{Caveat} For $i \ge |Fs|$ the subtraction truncates and selects the innermost binder;
        callers carry \texttt{TrProjCtor.field\_lt}.
  \end{itemize}
\end{definition}

\begin{definition}[Leading \texorpdfstring{$\Pi$}{Pi}-binder instantiation, \texttt{instPis}]
  \label{def:proj-inst-pis}
  \lean{Lean4Lean.VExpr.instPis}
  \leanok
  \uses{def:inst}
  \contrib{trproj}\srcloc{Lean4Lean/Theory/Proj.lean}{75}
  \thesisref{no direct counterpart; it models the parameter loop of the kernel's
  \texttt{inferProj}.}

  \lead{In short} \hl{Peels and instantiates} the leading $\Pi$-binders of a type, outermost
  first.
  \begin{itemize}
  \item $\mathrm{instPis}\;e\;[] = \mathrm{some}\;e$.
  \item $\mathrm{instPis}\;(\forall A.\,B)\;(a::as) = \mathrm{instPis}\;(B.\mathrm{inst}\;a)\;as$.
  \item Any other head with a non-empty argument list: $\mathrm{none}$.
  \item \lead{Caveat} Manifest binders only: a constructor type whose telescope hides behind a
        definition returns \texttt{none} (a completeness, not soundness, gap of
        \texttt{TrProjCtor.ctor}).
  \end{itemize}
\end{definition}

\begin{definition}[Telescope substitution at descending indices, \texttt{instFields}]
  \label{def:proj-inst-fields}
  \lean{Lean4Lean.VExpr.instFields}
  \leanok
  \uses{def:inst}
  \contrib{trproj}\srcloc{Lean4Lean/Theory/Proj.lean}{84}
  \thesisref{no counterpart; a de Bruijn convention of the formalisation.}

  \lead{In short} Substitutes a list of terms for the innermost binders of $F$, outermost
  argument first.
  \begin{itemize}
  \item $\mathrm{instFields}\;F\;[] = F$.
  \item $\mathrm{instFields}\;F\;(P::Ps) = \mathrm{instFields}\;(F.\mathrm{inst}\;P\;|Ps|)\;Ps$: the
        $m=|Ps|$ innermost binders $f_j=\mathrm{bvar}(m-1-j)$ are replaced by the $P_j$, each
        living over the outer context alone.
  \end{itemize}
\end{definition}

\begin{definition}[Motive body relative to given projections, \texttt{projMotiveBodyOf}]
  \label{def:proj-motive-body-of}
  \lean{Lean4Lean.VExpr.projMotiveBodyOf}
  \leanok
  \uses{def:proj-inst-fields, def:liftn}
  \contrib{trproj}\srcloc{Lean4Lean/Theory/Proj.lean}{91}
  \thesisref{Wtypes.tex \S5.1, the primitive $\Sigma$ rules $\pi_1\,p:\alpha$ and
  $\pi_2\,p:\beta[\pi_1\,p/x]$.}

  \lead{In short} The motive body for field $i$ once the earlier projections $Ps$ are fixed.
  \begin{itemize}
  \item $\mathrm{projMotiveBodyOf}\;Fs\;i\;Ps = \mathrm{instFields}\;F'\;(Ps.\mathrm{map}\;\lambda
        P.\,P.\mathrm{lift}\;(\mathrm{bvar}\;0))$, where $F'=(Fs.\mathrm{getD}\;i\;\mathrm{default}).
        \mathrm{liftN}\;1\;i$.
  \item Moves $F_i$ under the motive's new binder $x$, then replaces each earlier field $f_j$ by
        $P_j\,x$.
  \item \lead{Caveat} Out-of-range $i$ falls back to $\mathrm{default}=\mathrm{sort}\;\mathrm{zero}$:
        total but meaningless above the field count.
  \end{itemize}
\end{definition}

\begin{definition}[Projection function relative to given projections, \texttt{projFnOf}]
  \label{def:proj-fn-of}
  \lean{Lean4Lean.VExpr.projFnOf}
  \leanok
  \uses{def:proj-motive-body-of, def:proj-field-selector, def:mkapps}
  \contrib{trproj}\srcloc{Lean4Lean/Theory/Proj.lean}{96}
  \thesisref{typesys.tex \S3.1, the projection $\mathrm{inv}_x$ of $\mathrm{acc}\,x$ through
  $\mathrm{rec}_{\mathrm{acc}}$; axioms.tex \S2.6.3, the recursor.}

  \lead{In short} The recursor application computing field $i$ \hl{given} the earlier
  projections $Ps$.
  \begin{itemize}
  \item $\mathrm{projFnOf}\;S\;usS\;uss\;ps\;Fs\;i\;Ps$ is $\mathrm{const}\,(S.\mathrm{rec})\,(uss\,i)$
        applied to $ps \mathbin{+\!\!+} [\lambda x{:}(S\,usS)\,ps.\;\mathrm{projMotiveBodyOf}\;Fs\;i\;Ps,\;
        \mathrm{fieldSelector}\;Fs\;i]$.
  \item Separated from \ref{def:proj-fn} so the substitution lemmas below are stated once,
        generically in $Ps$.
  \end{itemize}
\end{definition}

\begin{definition}[The first \texorpdfstring{$i$}{i} projection functions, \texttt{projFns}]
  \label{def:proj-fns}
  \lean{Lean4Lean.VExpr.projFns}
  \leanok
  \uses{def:proj-fn-of}
  \contrib{trproj}\srcloc{Lean4Lean/Theory/Proj.lean}{108}
  \thesisref{axioms.tex \S2.6.3, the recursor: $u$ is a fresh universe variable per use of
  $\mathrm{rec}_P$.}

  \lead{In short} The list $[P_0,\dots,P_{i-1}]$, by recursion on $i$.
  \begin{itemize}
  \item $\mathrm{projFns}\dots 0 = []$.
  \item $\mathrm{projFns}\dots(i+1) = \mathrm{projFns}\dots i \mathbin{+\!\!+}
        [\mathrm{projFnOf}\dots i\;(\mathrm{projFns}\dots i)]$.
  \end{itemize}
\end{definition}

\begin{definition}[Motive body of projection \texorpdfstring{$i$}{i}, \texttt{projMotiveBody}]
  \label{def:proj-motive-body}
  \lean{Lean4Lean.VExpr.projMotiveBody}
  \leanok
  \uses{def:proj-motive-body-of, def:proj-fns}
  \contrib{trproj}\srcloc{Lean4Lean/Theory/Proj.lean}{114}
  \thesisref{Wtypes.tex \S5.1, $\pi_2\,p:\beta[\pi_1\,p/x]$.}

  \lead{In short} The codomain of $P_i$: the motive body once \emph{all} earlier projections are
  the real $P_j$'s.
  \begin{itemize}
  \item $\mathrm{projMotiveBody}\;S\;usS\;uss\;ps\;Fs\;i =
        \mathrm{projMotiveBodyOf}\;Fs\;i\;(\mathrm{projFns}\;S\;usS\;uss\;ps\;Fs\;i)$, over
        $\Gamma,x$.
  \item Used by \texttt{TrProjCtor.fn\_ty} (\ref{struct:trprojctor}): $P_i:\forall
        x{:}S\,usS\,ps,\;\mathrm{projMotiveBody}\dots i$.
  \end{itemize}
\end{definition}

\begin{definition}[The projection function, \texttt{projFn}]
  \label{def:proj-fn}
  \lean{Lean4Lean.VExpr.projFn}
  \leanok
  \uses{def:proj-motive-body, def:proj-field-selector, def:mkapps}
  \contrib{trproj}\srcloc{Lean4Lean/Theory/Proj.lean}{120}
  \thesisref{typesys.tex \S3.1, $\mathrm{inv}_x$ through $\mathrm{rec}_{\mathrm{acc}}$;
  axioms.tex \S2.6.3, the recursor.}

  \lead{In short} \hl{The representation of \texttt{Expr.proj} in the model}.
  \begin{itemize}
  \item $\mathrm{projFn}\;S\;usS\;uss\;ps\;Fs\;i$ is $\mathrm{const}\,(S.\mathrm{rec})\,(uss\,i)$
        applied to $ps \mathbin{+\!\!+} [\lambda x{:}(S\,usS)\,ps.\;\mathrm{projMotiveBody}\dots i,\;
        \mathrm{fieldSelector}\;Fs\;i]$.
  \item \texttt{TrProjCtor.eq} sets $e'=\mathrm{projFn}\dots i\;e$; definitionally
        $\mathrm{projFnOf}\dots i\;(\mathrm{projFns}\dots i)$ (\ref{fam:proj-unfolding}).
  \end{itemize}
\end{definition}

\begin{definition}[The type of a projected term, \texttt{projTy}]
  \label{def:proj-ty}
  \lean{Lean4Lean.VExpr.projTy}
  \leanok
  \uses{def:proj-inst-fields, def:proj-fns}
  \contrib{trproj}\srcloc{Lean4Lean/Theory/Proj.lean}{127}
  \thesisref{Wtypes.tex \S5.1, $\pi_2\,p:\beta[\pi_1\,p/x]$.}

  \lead{In short} The type the kernel's \texttt{inferProj} computes: each \texttt{.proj S j
  struct} replaced by $P_j\,struct$.
  \begin{itemize}
  \item $\mathrm{projTy}\;S\;usS\;uss\;ps\;Fs\;i\;e = \mathrm{instFields}\;
        (Fs.\mathrm{getD}\;i\;\mathrm{default})\;((\mathrm{projFns}\dots i).\mathrm{map}\;\lambda
        P.\,P\,e)$, over $\Gamma$.
  \item \lead{Caveat} Unused elsewhere except the example at
        \srcloc{Lean4Lean/Tests/ProjShape.lean}{178}; the identity relating it to
        \ref{def:proj-motive-body}, $\mathrm{projTy}\dots e =
        (\mathrm{projMotiveBody}\dots i).\mathrm{inst}\;e$, is asserted in its docstring but never
        proved, and the error message at \srcloc{Lean4Lean/Tests/ProjShape.lean}{114} still names
        it.
  \end{itemize}
\end{definition}

\begin{definition}[Arity of a telescope binder, \texttt{binderArity?}]
  \label{def:proj-binder-arity}
  \lean{Lean4Lean.VExpr.binderArity?}
  \leanok
  \uses{def:pi-telescope}
  \contrib{trproj}\srcloc{Lean4Lean/Theory/Proj.lean}{136}
  \thesisref{axioms.tex \S2.6.3: the recursor's minor premise arity.}

  \lead{In short} The $\Pi$-arity of the $k$-th binder of $ty$'s telescope.
  \begin{itemize}
  \item $ty.\mathrm{binderArity?}\;k = ty.\mathrm{piBinders}[k]?.\mathrm{map}\;\mathrm{piArity}$.
  \item As \texttt{TrProjCtor.minor\_arity}: $rci.\mathrm{type}.\mathrm{binderArity?}(np+1) =
        \mathrm{some}\;n_f$ says the recursor's minor binds exactly the fields, no inductive
        hypothesis.
  \item \lead{Caveat} Pins only the binder \emph{count}, not their types;
        \ref{thm:tr-env-proj-defeq} closes the type gap separately with kernel-side structure
        facts.
  \end{itemize}
\end{definition}

\begin{lemma}[Unfolding facts for the projection builders]
  \label{fam:proj-unfolding}
  \lean{Lean4Lean.VExpr.projFn_eq, Lean4Lean.VExpr.projFns_succ, Lean4Lean.VExpr.projFns_length, Lean4Lean.VExpr.projMotiveBody_zero}
  \leanok
  \uses{def:proj-fn, def:proj-fns, def:proj-motive-body}
  \contrib{trproj}\srcloc{Lean4Lean/Theory/Proj.lean}{140}
  \thesisref{no counterpart; definitional bookkeeping.}

  \lead{In short} Four bookkeeping identities relating $\mathrm{projFn}$, $\mathrm{projFns}$ and
  $\mathrm{projMotiveBody}$.
  \begin{itemize}
  \item $\mathrm{projFn}\dots i = \mathrm{projFnOf}\dots i\;(\mathrm{projFns}\dots i)$.
  \item $\mathrm{projFns}\dots(i+1) = \mathrm{projFns}\dots i \mathbin{+\!\!+}
        [\mathrm{projFn}\dots i]$.
  \item $(\mathrm{projFns}\dots i).\mathrm{length} = i$.
  \item The motive of field $0$ is the constant motive $(Fs.\mathrm{getD}\;0\;
        \mathrm{default}).\mathrm{lift}$.
  \end{itemize}
\end{lemma}
\begin{proof}
  \leanok
  \lead{Idea} All four are \texttt{rfl} except \texttt{projFns\_length}, an induction on $i$ via
  \texttt{projFns\_succ}.
\end{proof}

\section{The builders under substitution}

Every builder is proved once against \texttt{Subst} (\ref{def:subst}), shifting by one binder per
field, and weakening/instantiation are read off; level instantiation keeps its own \texttt{instL}
proofs.

\begin{lemma}[Lambda-telescopes under substitution]
  \label{fam:proj-foldr-lam}
  \lean{Lean4Lean.VExpr.foldr_lam_subst, Lean4Lean.VExpr.foldr_lam_inst, Lean4Lean.VExpr.foldr_lam_instL}
  \leanok
  \uses{def:subst}
  \contrib{trproj}\srcloc{Lean4Lean/Theory/Proj.lean}{164}
  \thesisref{no counterpart; de Bruijn bookkeeping.}

  \lead{In short} Pushing a substitution through a $\lambda$-telescope.
  \begin{itemize}
  \item $(\mathrm{foldr}\;\lambda\;base\;Fs).\mathrm{subst}\;\sigma = \mathrm{foldr}\;\lambda\;
        (base.\mathrm{subst}\;(\sigma.\mathrm{liftN}\;|Fs|))\;(Fs.\mathrm{mapIdx}\;\lambda j\,F.\,
        F.\mathrm{subst}\;(\sigma.\mathrm{liftN}\;j))$.
  \item The instantiation form has $\mathrm{inst}\;e_0\;(k+j)$ in place of $\mathrm{subst}\;\sigma$;
        the level form maps $\mathrm{instL}$ over the telescope.
  \end{itemize}
\end{lemma}
\begin{proof}
  \uses{fam:subst-liftn-algebra, fam:subst-basic}
  \leanok
  \lead{Idea} \texttt{foldr\_lam\_subst}: structural recursion via \texttt{Subst.lift\_liftN}.
  \texttt{foldr\_lam\_inst}, \texttt{foldr\_lam\_instL} follow.
\end{proof}

\begin{lemma}[The field selector under substitution]
  \label{fam:proj-fieldselector-subst}
  \lean{Lean4Lean.VExpr.fieldSelector_subst, Lean4Lean.VExpr.fieldSelector_instL}
  \leanok
  \uses{def:proj-field-selector, def:subst}
  \contrib{trproj}\srcloc{Lean4Lean/Theory/Proj.lean}{185}
  \thesisref{no counterpart; de Bruijn bookkeeping.}

  \lead{In short} Substitution distributes over $\mathrm{fieldSelector}$, for $i<|Fs|$.
  \begin{itemize}
  \item $(\mathrm{fieldSelector}\;Fs\;i).\mathrm{subst}\;\sigma =
        \mathrm{fieldSelector}\;(Fs.\mathrm{mapIdx}\;\lambda j\,F.\,F.\mathrm{subst}\;
        (\sigma.\mathrm{liftN}\;j))\;i$.
  \item The level-instantiation form holds with no side condition.
  \end{itemize}
\end{lemma}
\begin{proof}
  \uses{fam:proj-foldr-lam, fam:subst-liftn-algebra}
  \leanok
  \lead{Idea} Rewrite by \ref{fam:proj-foldr-lam} and \texttt{subst\_bvar}, then
  \texttt{Subst.Fixes.liftN} at the selected variable $|Fs|-1-i$, where $i<|Fs|$ is used; the
  \texttt{instL} form follows by \texttt{simp}.
\end{proof}

\begin{lemma}[Equations for \texttt{instFields}]
  \label{fam:proj-instfields-simp}
  \lean{Lean4Lean.VExpr.instFields_nil, Lean4Lean.VExpr.instFields_cons}
  \leanok
  \uses{def:proj-inst-fields}
  \contrib{trproj}\srcloc{Lean4Lean/Theory/Proj.lean}{205}
  \thesisref{no counterpart; definitional bookkeeping.}

  \lead{In short} The two defining equations of $\mathrm{instFields}$, restated as rewrite rules.
  \begin{itemize}
  \item $\mathrm{instFields}\;F\;[]=F$ (\texttt{@[simp]}).
  \item $\mathrm{instFields}\;F\;(P::Ps)=\mathrm{instFields}\;(F.\mathrm{inst}\;P\;|Ps|)\;Ps$.
  \end{itemize}
\end{lemma}
\begin{proof}
  \leanok
  \lead{Idea} Both \texttt{rfl}.
\end{proof}

\begin{lemma}[\texttt{instFields} under substitution and level instantiation]
  \label{fam:proj-instfields-subst}
  \lean{Lean4Lean.VExpr.instFields_subst, Lean4Lean.VExpr.instFields_instL}
  \leanok
  \uses{def:proj-inst-fields, def:subst}
  \contrib{trproj}\srcloc{Lean4Lean/Theory/Proj.lean}{210}
  \thesisref{no counterpart; de Bruijn bookkeeping.}

  \lead{In short} Substitution and level instantiation distribute over $\mathrm{instFields}$.
  \begin{itemize}
  \item $(\mathrm{instFields}\;F\;Ps).\mathrm{subst}\;\sigma = \mathrm{instFields}\;
        (F.\mathrm{subst}\;(\sigma.\mathrm{liftN}\;|Ps|))\;(Ps.\mathrm{map}\;(\cdot.\mathrm{subst}\;
        \sigma))$.
  \item Same with $\mathrm{instL}\;ls$, with no lifting (level instantiation moves no variables).
  \end{itemize}
\end{lemma}
\begin{proof}
  \uses{thm:subst-instn, fam:inst-inst, fam:proj-instfields-simp}
  \leanok
  \lead{Idea} Structural recursion on $Ps$, closing each step with \texttt{subst\_instN}
  (resp.\ \texttt{instL\_instN}).
\end{proof}

\begin{lemma}[\texttt{instPis} commutes with substitution]
  \label{fam:proj-instpis-subst}
  \lean{Lean4Lean.VExpr.instPis_subst, Lean4Lean.VExpr.instPis_lift', Lean4Lean.VExpr.instPis_inst, Lean4Lean.VExpr.instPis_instL}
  \leanok
  \uses{def:proj-inst-pis, def:subst}
  \contrib{trproj}\srcloc{Lean4Lean/Theory/Proj.lean}{226}
  \thesisref{no counterpart; de Bruijn bookkeeping.}

  \lead{In short} \hl{Success is preserved}, not merely the value, under substitution.
  \begin{itemize}
  \item If $T.\mathrm{instPis}\;ps=\mathrm{some}\;cty$ then $(T.\mathrm{subst}\;\sigma).
        \mathrm{instPis}\;(ps.\mathrm{map}\;(\cdot.\mathrm{subst}\;\sigma)) =
        \mathrm{some}\;(cty.\mathrm{subst}\;\sigma)$.
  \item Likewise for $\mathrm{lift'}\;\rho$, $\mathrm{inst}\;e_0\;k$ and $\mathrm{instL}\;ls$.
  \end{itemize}
\end{lemma}
\begin{proof}
  \uses{fam:subst-basic, thm:lift-eq-subst}
  \leanok
  \lead{Idea} Recursion on $ps$ with \texttt{subst\_inst} at each $\Pi$-binder, \texttt{nomatch} on
  the five non-$\Pi$ heads; the \texttt{lift'} and \texttt{inst} forms are \texttt{simpa}
  corollaries.
\end{proof}

\begin{lemma}[\texttt{instPis} preserves \texttt{CtorHeaded}]
  \label{thm:proj-instpis-ctorheaded}
  \lean{Lean4Lean.VExpr.instPis_ctorHeaded}
  \leanok
  \uses{def:proj-inst-pis, def:ctorheaded}
  \contrib{trproj}\srcloc{Lean4Lean/Theory/Proj.lean}{254}
  \thesisref{no counterpart; a stability property of the formalisation's shape predicates.}

  \lead{In short} $\mathrm{instPis}$ keeps a constructor-headed type constructor-headed.
  \begin{itemize}
  \item If $T$ is \texttt{CtorHeaded} and $T.\mathrm{instPis}\;ps=\mathrm{some}\;cty$, then $cty$
        is \texttt{CtorHeaded}.
  \end{itemize}
\end{lemma}
\begin{proof}
  \uses{fam:ctorheaded-stability}
  \leanok
  \lead{Idea} Recursion on $ps$, applying \texttt{CtorHeaded.inst} at each $\Pi$-binder; the
  non-$\Pi$ heads cannot occur.
\end{proof}

\begin{lemma}[\texttt{instPis} consumes exactly its arguments]
  \label{thm:proj-instpis-piarity}
  \lean{Lean4Lean.VExpr.instPis_piArity}
  \leanok
  \uses{def:proj-inst-pis, def:ctorheaded, def:pi-telescope}
  \contrib{trproj}\srcloc{Lean4Lean/Theory/Proj.lean}{264}
  \thesisref{axioms.tex \S2.6.1, inductive specifications: a constructor's arity.}

  \lead{In short} Consumed arguments plus remaining arity equal the original arity.
  \begin{itemize}
  \item For \texttt{CtorHeaded} $T$ with $T.\mathrm{instPis}\;ps=\mathrm{some}\;cty$:
        $cty.\mathrm{piArity}+|ps| = T.\mathrm{piArity}$.
  \item Turns a constructor's kernel arity into the model's $np+|fieldTys|$ in
        \ref{thm:tr-env-proj-defeq}.
  \end{itemize}
\end{lemma}
\begin{proof}
  \uses{fam:ctorheaded-stability}
  \leanok
  \lead{Idea} Recursion on $ps$ via \texttt{piArity\_inst\_of\_ctorHeaded} (instantiation creates
  no binders), closed by \texttt{omega}.
\end{proof}

\begin{lemma}[The motive body under substitution]
  \label{fam:proj-motivebodyof-subst}
  \lean{Lean4Lean.VExpr.projMotiveBodyOf_subst, Lean4Lean.VExpr.projMotiveBodyOf_instL}
  \leanok
  \uses{def:proj-motive-body-of, def:subst}
  \contrib{trproj}\srcloc{Lean4Lean/Theory/Proj.lean}{276}
  \thesisref{no counterpart; de Bruijn bookkeeping.}

  \lead{In short} Substitution distributes over $\mathrm{projMotiveBodyOf}$, shifted by the
  motive's own binder.
  \begin{itemize}
  \item Given $|Ps|=i$: $(\mathrm{projMotiveBodyOf}\;Fs\;i\;Ps).\mathrm{subst}\;\sigma.\mathrm{lift}$
        is the motive body of the substituted field telescope (binder-wise) and the substituted
        earlier projections (pointwise).
  \item The $\mathrm{instL}$ form holds unconditionally.
  \end{itemize}
\end{lemma}
\begin{proof}
  \uses{fam:proj-instfields-subst, thm:liftn-subst-liftn, fam:subst-liftn-algebra}
  \leanok
  \lead{Idea} Unfold and rewrite by \ref{fam:proj-instfields-subst}, \texttt{Subst.lift\_liftN},
  \texttt{liftN\_subst\_liftN} and \texttt{List.getD\_mapIdx}; close the congruence on $P_j\,x$
  with \texttt{lift\_subst\_lift}.
\end{proof}

\begin{lemma}[The parameterised projection function under substitution]
  \label{fam:proj-fnof-subst}
  \lean{Lean4Lean.VExpr.projFnOf_subst, Lean4Lean.VExpr.projFnOf_instL}
  \leanok
  \uses{def:proj-fn-of, def:subst}
  \contrib{trproj}\srcloc{Lean4Lean/Theory/Proj.lean}{294}
  \thesisref{no counterpart; de Bruijn bookkeeping.}

  \lead{In short} Substitution distributes over $\mathrm{projFnOf}$, needing $i<|Fs|$ and $|Ps|=i$.
  \begin{itemize}
  \item Pushes into the parameters, the field telescope (binder-wise) and the earlier projections.
  \item Level instantiation distributes unconditionally: $\mathrm{VLevel.inst}\;ls$ over $usS$, and
        pointwise via $\lambda j.\,(uss\,j).\mathrm{map}\;(\mathrm{VLevel.inst}\;ls)$.
  \end{itemize}
\end{lemma}
\begin{proof}
  \uses{fam:mkapps-stability, fam:proj-motivebodyof-subst, fam:proj-fieldselector-subst}
  \leanok
  \lead{Idea} Rewrite by \texttt{mkApps\_subst} (resp.\ \texttt{mkApps\_instL}), then
  \ref{fam:proj-motivebodyof-subst} and \ref{fam:proj-fieldselector-subst} for the two trailing
  arguments.
\end{proof}

\begin{lemma}[The list of projection functions under substitution]
  \label{fam:proj-fns-subst}
  \lean{Lean4Lean.VExpr.projFns_subst, Lean4Lean.VExpr.projFns_instL}
  \leanok
  \uses{def:proj-fns, def:subst}
  \contrib{trproj}\srcloc{Lean4Lean/Theory/Proj.lean}{313}
  \thesisref{no counterpart; de Bruijn bookkeeping.}

  \lead{In short} Mapping a substitution over $\mathrm{projFns}\dots i$ gives the projection
  functions of the substituted data, for $i\le|Fs|$.
  \begin{itemize}
  \item The level-instantiation form has no side condition.
  \end{itemize}
\end{lemma}
\begin{proof}
  \uses{fam:proj-fnof-subst, fam:proj-unfolding}
  \leanok
  \lead{Idea} Induction on $i$ via \texttt{projFns\_succ}, applying \ref{fam:proj-fnof-subst} and
  the induction hypothesis \emph{twice}, since \ref{def:proj-fns}'s defining equation mentions the
  predecessor list twice.
\end{proof}

\begin{lemma}[The projection function under substitution]
  \label{fam:proj-fn-subst}
  \lean{Lean4Lean.VExpr.projFn_subst, Lean4Lean.VExpr.projFn_lift', Lean4Lean.VExpr.projFn_inst, Lean4Lean.VExpr.projFn_instL}
  \leanok
  \uses{def:proj-fn, def:subst}
  \contrib{trproj}\srcloc{Lean4Lean/Theory/Proj.lean}{333}
  \thesisref{no counterpart; de Bruijn bookkeeping.}

  \lead{In short} $\mathrm{projFn}$ \hl{commutes} with $\mathrm{subst}$, $\mathrm{lift'}$,
  $\mathrm{inst}$ and $\mathrm{instL}$.
  \begin{itemize}
  \item Pushes into the parameters and, binder-wise, into the field telescope.
  \item The three term-level forms carry the side condition $i<|Fs|$.
  \end{itemize}
\end{lemma}
\begin{proof}
  \uses{fam:proj-fnof-subst, fam:proj-fns-subst, fam:proj-unfolding, thm:lift-eq-subst}
  \leanok
  \lead{Idea} Rewrite by \texttt{projFn\_eq} into $\mathrm{projFnOf}$, apply
  \ref{fam:proj-fnof-subst} and \ref{fam:proj-fns-subst}; the \texttt{lift'}/\texttt{inst} forms
  follow via \texttt{lift'\_eq\_subst} and \texttt{instN\_eq}.
\end{proof}

\begin{lemma}[The motive body under weakening and instantiation]
  \label{fam:proj-motivebody-subst}
  \lean{Lean4Lean.VExpr.projMotiveBody_subst, Lean4Lean.VExpr.projMotiveBody_lift', Lean4Lean.VExpr.projMotiveBody_instN, Lean4Lean.VExpr.projMotiveBody_instL}
  \leanok
  \uses{def:proj-motive-body, def:subst}
  \contrib{trproj}\srcloc{Lean4Lean/Theory/Proj.lean}{357}
  \thesisref{no counterpart; de Bruijn bookkeeping.}

  \lead{In short} The same four forms for $\mathrm{projMotiveBody}$, shifted by the motive's own
  binder $x$.
  \begin{itemize}
  \item $\sigma.\mathrm{lift}$ for \texttt{subst}, $\rho.\mathrm{cons}$ for \texttt{lift'},
        $\mathrm{inst}\;e_0\;(k+1)$ for \texttt{instN}.
  \item Term-level forms carry $i\le|Fs|$; the $\mathrm{instL}$ form carries none.
  \end{itemize}
\end{lemma}
\begin{proof}
  \uses{fam:proj-motivebodyof-subst, fam:proj-fns-subst, fam:subst-liftn-algebra}
  \leanok
  \lead{Idea} Unfold to $\mathrm{projMotiveBodyOf}$, apply \ref{fam:proj-motivebodyof-subst} and
  \ref{fam:proj-fns-subst}; the remaining forms are \texttt{Subst} bookkeeping.
\end{proof}

\section{Field selection on variable spines}

Saturating a $\lambda$-telescope by its arguments substitutes them outermost first, exactly
$\mathrm{instFields}$ -- what \ref{thm:tr-env-proj-defeq} needs for the iota reduct's template.

\begin{lemma}[\texttt{instFields} on applications and lifted terms]
  \label{fam:proj-instfields-spine}
  \lean{Lean4Lean.VExpr.instFields_app, Lean4Lean.VExpr.instFields_mkApps, Lean4Lean.VExpr.instFields_liftN}
  \leanok
  \uses{def:proj-inst-fields, def:mkapps, def:liftn}
  \contrib{trproj}\srcloc{Lean4Lean/Theory/Proj.lean}{394}
  \thesisref{no counterpart; de Bruijn bookkeeping.}

  \lead{In short} $\mathrm{instFields}$ distributes over application, and cancels a matching lift.
  \begin{itemize}
  \item Distributes over $\mathrm{app}$ and over $\mathrm{mkApps}$.
  \item $\mathrm{instFields}\;(e.\mathrm{liftN}\;|Ps|)\;Ps = e$.
  \end{itemize}
\end{lemma}
\begin{proof}
  \uses{fam:proj-instfields-simp, def:inst}
  \leanok
  \lead{Idea} Structural recursion on $Ps$; \texttt{instFields\_mkApps} recurses on the argument
  list, using \texttt{liftN'\_liftN\_lo} and \texttt{inst\_liftN} to cancel one lift per step.
\end{proof}

\begin{lemma}[\texttt{instFields} on a bound variable]
  \label{thm:proj-instfields-bvar}
  \lean{Lean4Lean.VExpr.instFields_bvar}
  \leanok
  \uses{def:proj-inst-fields}
  \contrib{trproj}\srcloc{Lean4Lean/Theory/Proj.lean}{413}
  \thesisref{no counterpart; de Bruijn bookkeeping.}

  \lead{In short} A selected bound variable becomes the correspondingly indexed argument.
  \begin{itemize}
  \item For $j<|Ps|$: $\mathrm{instFields}\;(\mathrm{bvar}\;j)\;Ps = Ps.\mathrm{getD}\;
        (|Ps|-1-j)\;\mathrm{default}$.
  \end{itemize}
\end{lemma}
\begin{proof}
  \uses{fam:proj-instfields-spine}
  \leanok
  \lead{Idea} Recursion on $Ps$, splitting \texttt{instVar} into its three cases; the hit case
  closes with \texttt{instFields\_liftN} and \texttt{omega}.
\end{proof}

\begin{theorem}[The iota reduct on the minor and field spine]
  \label{thm:proj-instfields-minor-spine}
  \lean{Lean4Lean.VExpr.instFields_minor_spine}
  \leanok
  \uses{def:proj-inst-fields, def:mkapps, def:bvarsdesc}
  \contrib{trproj}\srcloc{Lean4Lean/Theory/Proj.lean}{428}
  \thesisref{axioms.tex \S2.6.4, the computation rule (iota reduction):
  $\mathrm{rec}_P\,C\,e\,p[b]\,(c\,b)\equiv e_c\,b\,v$.}

  \lead{In short} Saturating the iota template selects the minor premise and applies it to the
  fields.
  \begin{itemize}
  \item $\mathrm{instFields}\;((\mathrm{bvar}\;|fs|).\mathrm{mkApps}\;
        (\mathrm{bvarsDesc}\;0\;|fs|))\;(pre \mathbin{+\!\!+} s :: fs) = s.\mathrm{mkApps}\;fs$: the
        template $\lambda\,\mathit{params}\,\mathit{motive}\,\mathit{minor}\,\mathit{fields}.\,
        \mathit{minor}\;\mathit{fields}$, saturated.
  \end{itemize}
\end{theorem}
\begin{proof}
  \uses{fam:proj-instfields-spine, thm:proj-instfields-bvar}
  \leanok
  \lead{Idea} Rewrite by \texttt{instFields\_mkApps}, then \ref{thm:proj-instfields-bvar} on the
  head $\mathrm{bvar}\;|fs|$ (landing on $s$ since $|pre \mathbin{+\!\!+} s::fs|-1-|fs|=|pre|$) and on
  each argument via \texttt{getElem\_bvarsDesc} and \texttt{omega}. The arithmetic core of
  \ref{thm:tr-env-proj-defeq}'s $\beta$ phase.
\end{proof}

\section{Deciding the syntactic predicates}

\texttt{Tests/ShapeDecide.lean} makes every syntactic predicate of Chapter~\ref{chap:inductive}
executable, so the tests below can \texttt{decide} them on the kernel's own data.

\begin{example}[Decidable equality for \texttt{VLevel} and \texttt{VExpr}]
  \label{fam:shapedecide-deceq}
  \lean{Lean4Lean.instDecidableEqVLevel, Lean4Lean.instDecidableEqVExpr}
  \leanok
  \uses{ind:vexpr, ind:vlevel}
  \contrib{iota}\srcloc{Lean4Lean/Tests/ShapeDecide.lean}{15}
  \thesisref{no counterpart; executability scaffolding.}

  \lead{In short} \texttt{deriving DecidableEq} for \texttt{VLevel} and \texttt{VExpr}.
  \begin{itemize}
  \item Makes every \texttt{decide}/\texttt{unless} in the shape tests, and every \texttt{VExpr}
        comparison in \srcloc{Lean4Lean/Tests/IotaShape.lean}{1},
        \srcloc{Lean4Lean/Tests/ProjShape.lean}{1} and \srcloc{Lean4Lean/Tests/ProjInhabit.lean}{1},
        run.
  \end{itemize}
\end{example}

\begin{example}[Decidability of the existential idioms]
  \label{fam:shapedecide-existential}
  \lean{Lean4Lean.VExpr.instDecidableExistsListEqMkAppsHAppend, Lean4Lean.VExpr.instDecidableExistsListAndEqNatMkAppsHAppend, Lean4Lean.VExpr.instDecidableExistsAndEqOptionSomeOfDecidablePred}
  \leanok
  \uses{thm:eq-mkapps-append, def:mkapps}
  \contrib{iota}\srcloc{Lean4Lean/Tests/ShapeDecide.lean}{20}
  \thesisref{no counterpart; executability scaffolding.}

  \lead{In short} Three instances making recurring existential shapes executable.
  \begin{itemize}
  \item $\exists\,rest,\;e=f.\mathrm{mkApps}\;(pre \mathbin{+\!\!+} rest)$, and the same with a
        length constraint on $rest$: \texttt{decidable\_of\_iff} against
        \texttt{eq\_mkApps\_append\_iff}.
  \item $\exists a,\;o=\mathrm{some}\;a\wedge P\,a$ for an \texttt{Option} and decidable $P$: a
        case split on the option.
  \end{itemize}
\end{example}

\begin{example}[Decidability of the head-constructor tests]
  \label{fam:shapedecide-head}
  \lean{Lean4Lean.VExpr.instDecidableExistsVLevelEqSort, Lean4Lean.VExpr.instDecidableExistsNatEqBvar, Lean4Lean.VExpr.instDecidableExistsNameListVLevelEqConst, Lean4Lean.VExpr.instDecidableRecHeaded, Lean4Lean.VExpr.instDecidableCtorHeaded}
  \leanok
  \uses{def:head-tests, def:recheaded, def:ctorheaded}
  \contrib{trproj}\srcloc{Lean4Lean/Tests/ShapeDecide.lean}{34}
  \thesisref{no counterpart; executability scaffolding.}

  \lead{In short} Decidability for $\exists u,e{=}\mathrm{sort}\;u$, $\exists k,e{=}\mathrm{bvar}\;
  k$, $\exists I\,us,e{=}\mathrm{const}\;I\;us$, and for \texttt{RecHeaded}/\texttt{CtorHeaded}.
  \begin{itemize}
  \item Each by \texttt{decidable\_of\_iff} against \texttt{isSort\_iff}/\texttt{isBvar\_iff}/
        \texttt{isConst\_iff}.
  \item \texttt{RecHeaded} and \texttt{CtorHeaded} apply the characterisation to
        $ty.\mathrm{piBody}.\mathrm{getAppFn}$.
  \end{itemize}
\end{example}

\begin{example}[Decision procedures for the block shape predicates]
  \label{fam:shapedecide-shapes}
  \lean{Lean4Lean.VExpr.instDecidableMentionsConst, Lean4Lean.VExpr.instDecidableCtorResult, Lean4Lean.VExpr.instDecidableMajorApp, Lean4Lean.VExpr.instDecidableValidIndApp, Lean4Lean.VExpr.instDecidableFieldPositive, Lean4Lean.VExpr.instDecidableCtorPositive, Lean4Lean.VExpr.instDecidableFieldInIndices, Lean4Lean.VExpr.instDecidableMotiveShape, Lean4Lean.VExpr.instDecidableMinorHeaded, Lean4Lean.VExpr.instDecidableMinorFor, Lean4Lean.VExpr.instDecidableRecShape, Lean4Lean.VExpr.instDecidableCtorShape, Lean4Lean.VExpr.instDecidableRuleShape, Lean4Lean.instDecidableLargeElimShape}
  \leanok
  \uses{family:ind-shape-iff, def:ind-mentions-const, def:ind-large-elim-shape, structure:ind-wf}
  \contrib{iota}\srcloc{Lean4Lean/Tests/ShapeDecide.lean}{41}
  \thesisref{axioms.tex \S2.6.1, inductive specifications; axioms.tex \S2.6.3, the recursor.}

  \lead{In short} Executable instances for every syntactic predicate of Chapter~\ref{chap:inductive},
  from \texttt{MentionsConst} and \texttt{CtorPositive} to \texttt{RuleShape} and
  \texttt{VInductDecl.LargeElimShape}.
  \begin{itemize}
  \item Four go through hand-written \texttt{\_iff} lemmas; the rest are
        \texttt{unfold; infer\_instance}.
  \item \lead{Caveat} All twenty-two instances are global \texttt{Decidable} instances declared
        from a \texttt{Tests} module; safe only while no library module imports it. The widest,
        \srcloc{Lean4Lean/Tests/ShapeDecide.lean}{27}, decides an arbitrary \texttt{Option}
        existential.
  \end{itemize}
\end{example}

\section{Validating the inductive specification against the kernel}

\texttt{Tests/IotaShape.lean} decides every syntactic clause of \texttt{VInductDecl.WF}
(\ref{structure:ind-wf}) on the kernel's own recursor, constructor and rule data, and compares the
model's iota reduct builder with the kernel's.

\begin{example}[Direct-block fixtures and the documented nested control]
  \label{def:iotashape-fixtures-direct}
  \lean{Lean4Lean.Tests.IotaShape.P2, Lean4Lean.Tests.IotaShape.Vec, Lean4Lean.Tests.IotaShape.Ev, Lean4Lean.Tests.IotaShape.Od, Lean4Lean.Tests.IotaShape.Tree}
  \leanok
  \contrib{iota}\srcloc{Lean4Lean/Tests/IotaShape.lean}{39}
  \thesisref{axioms.tex \S2.6.1, inductive specifications.}

  \lead{In short} Four local declarations exercising \ref{structure:ind-wf}'s direct-block
  clauses.
  \begin{itemize}
  \item \texttt{P2}: a two-field \texttt{Nat} structure. \texttt{Vec}: indexed over
        $\mathrm{Nat}$. \texttt{Ev}/\texttt{Od}: a mutual \texttt{Prop} block.
  \item \texttt{Tree} ($\mathrm{node}:\mathrm{List}\;\mathrm{Tree}\to\mathrm{Tree}$): the negative
        control. Lean accepts it as \emph{nested}, but its constructor re-enters the block through
        \texttt{List}, so \texttt{ctors\_positive} (\ref{def:ind-ctor-positive}) rejects it
        directly.
  \end{itemize}
\end{example}

\begin{example}[The \texttt{rule\_shape} clause for one rule, \texttt{ruleShapeAt}]
  \label{def:iotashape-ruleshapeat}
  \lean{Lean4Lean.Tests.IotaShape.ruleShapeAt, Lean4Lean.Tests.IotaShape.instDecidableRuleShapeAt}
  \leanok
  \uses{def:ind-rule-shape, def:ind-minor-for, def:pi-telescope, structure:ind-wf}
  \contrib{iota}\srcloc{Lean4Lean/Tests/IotaShape.lean}{59}
  \thesisref{axioms.tex \S2.6.4, the computation rule (iota reduction); axioms.tex \S2.6.3, the
  recursor.}

  \lead{In short} Asserts a minor premise for constructor $c$ whose reduct has the recursor's
  \texttt{RuleShape}.
  \begin{itemize}
  \item $\mathrm{ruleShapeAt}\;ty\;rhs\;np\;nm\;nmin\;nf\;c$: some minor index $j<nmin$ at
        telescope position $np+nm+j$ is \texttt{MinorFor} $c$, has at least $nf$ binders, and
        $rhs$ satisfies $\mathtt{RuleShape}\;np\;nm\;nmin\;nf\;(A.\mathrm{piArity}-nf)\;j$.
  \item Mirrors the \texttt{rule\_shape} clause of \ref{structure:ind-wf}; decidable by
        \texttt{unfold; infer\_instance}.
  \end{itemize}
\end{example}

\begin{example}[Deciding the per-recursor clauses, \texttt{checkShapes}]
  \label{def:iotashape-checkshapes}
  \lean{Lean4Lean.Tests.IotaShape.checkShapes}
  \leanok
  \uses{def:iotashape-ruleshapeat, def:ind-rec-shape, def:ind-ctor-shape, def:recheaded, def:ctorheaded, def:ind-major-former, def:meta-ofexpr}
  \contrib{iota}\srcloc{Lean4Lean/Tests/IotaShape.lean}{72}
  \thesisref{axioms.tex \S2.6.3, the recursor.}

  \lead{In short} Decides a kernel recursor's shape clauses of \ref{structure:ind-wf} at its own
  telescope split.
  \begin{itemize}
  \item \texttt{RecShape}, \texttt{RecHeaded}, failure of \texttt{CtorHeaded}, and
        \texttt{majorFormer?}; then \texttt{recs\_elim} (the recursor's universe count is the
        block's or exactly one more, matched by each motive ending in $\mathrm{sort}\,\mathrm{zero}$
        or $\mathrm{sort}\,(\mathrm{param}\;0)$) and \texttt{rules\_total}.
  \item Per rule (\ref{def:iotashape-ruleshapeat}): \texttt{CtorShape}, failure of
        \texttt{RecHeaded}, and agreement of \texttt{nfields} with \texttt{numFields}.
  \item \lead{Caveat} \texttt{recs\_elim} assumes the fresh universe sits at index $0$ of the
        recursor's \texttt{levelParams}, unchecked here.
  \end{itemize}
\end{example}

\begin{example}[Executable pattern matcher, \texttt{matchPat}]
  \label{def:iotashape-matchpat}
  \lean{Lean4Lean.Tests.IotaShape.matchPat}
  \leanok
  \uses{inductive:pattern-core, inductive:pattern-matches, ind:vexpr}
  \contrib{iota}\srcloc{Lean4Lean/Tests/IotaShape.lean}{103}
  \thesisref{axioms.tex \S2.6.4: the computation rule.}

  \lead{In short} Executable pattern matching, returning a level and argument assignment.
  \begin{itemize}
  \item $\mathrm{const}\;c$ matches $\mathrm{const}\;c'\;ls$ when $c=c'$: returns $ls$, the empty
        assignment.
  \item $\mathrm{var}\;f$ matches an application: records the argument at the new path, recurses
        into the function.
  \item $\mathrm{app}\;f\;a$: recurses into both halves, joins with $\mathrm{Sum.elim}$.
  \item Any other combination: $\mathrm{none}$.
  \end{itemize}
\end{example}

\begin{example}[Soundness of the executable matcher]
  \label{thm:iotashape-matchpat-sound}
  \lean{Lean4Lean.Tests.IotaShape.matchPat_sound}
  \leanok
  \uses{def:iotashape-matchpat, inductive:pattern-matches}
  \contrib{iota}\srcloc{Lean4Lean/Tests/IotaShape.lean}{115}
  \thesisref{axioms.tex \S2.6.4: the computation rule.}

  \lead{In short} The matcher only ever reports \hl{genuine matches}.
  \begin{itemize}
  \item If $\mathrm{matchPat}\;p\;e=\mathrm{some}\,(m_1,m_2)$ then $p.\mathrm{Matches}\;e\;m_1\;
        m_2$.
  \item \lead{Idea} Recursion on pattern and expression together; the three matching cases close
        via \texttt{simp}/\texttt{split}, the fifteen impossible shapes in one \texttt{simp}.
  \item \lead{Caveat} Completeness is not proved: a failed match errors, the right asymmetry for a
        test.
  \end{itemize}
\end{example}

\begin{example}[Comparing the iota reduct with the kernel's]
  \label{def:iotashape-checkiota}
  \lean{Lean4Lean.Tests.IotaShape.checkIota, Lean4Lean.Tests.IotaShape.checkIotaAuto}
  \leanok
  \uses{def:pattern-iota-rhs, inductive:simple-pattern, def:iotashape-matchpat, def:kernel-inductive-reduce-rec, def:meta-ofexpr, def:closedn}
  \contrib{iota}\srcloc{Lean4Lean/Tests/IotaShape.lean}{141}
  \thesisref{axioms.tex \S2.6.4: the computation rule.}

  \lead{In short} Compares the model's iota-reduct builder with the kernel's own reduction.
  \begin{itemize}
  \item \texttt{checkIota}: reduces the redex $rec\,ps\,Ms\,ms\,is\,(ctor\,cps\,fs)$ with
        \texttt{inductiveReduceRec}, translates both sides to \texttt{VExpr}, and requires
        $\mathtt{iotaRHS}\dots.\mathrm{apply}$ at the matched pattern to equal the kernel's reduct
        \emph{syntactically} (the stored template must be \texttt{Closed}).
  \item \texttt{checkIotaAuto}: reads levels and parameters off the major premise, covering
        nested-block auxiliary recursors generically.
  \item \lead{Caveat} The oracle is lean4lean's own \texttt{inductiveReduceRec}
        (\srcloc{Lean4Lean/Inductive/Reduce.lean}{100}), not the C++ kernel. The guard at
        \srcloc{Lean4Lean/Tests/IotaShape.lean}{306} is a \texttt{trproj} edit.
  \end{itemize}
\end{example}

\begin{example}[The remaining per-recursor clauses]
  \label{def:iotashape-checkmore}
  \lean{Lean4Lean.Tests.IotaShape.checkTypesHaveRec, Lean4Lean.Tests.IotaShape.checkMore}
  \leanok
  \uses{def:ind-major-former, def:ind-ctor-shape, def:closedn, def:meta-ofexpr}
  \contrib{iota}\srcloc{Lean4Lean/Tests/IotaShape.lean}{171}
  \thesisref{axioms.tex \S2.6.3, the recursor.}

  \lead{In short} Decides the clauses \ref{def:iotashape-checkshapes} leaves out.
  \begin{itemize}
  \item \texttt{checkTypesHaveRec}: $I$'s recursor eliminates $I$ itself, via
        \texttt{majorFormer?}.
  \item \texttt{checkMore}: \texttt{rec\_params}, \texttt{rules\_nodup},
        \texttt{rules\_own\_params}, each rule's constructor \texttt{CtorShape}, and closedness of
        every rule reduct.
  \end{itemize}
\end{example}

\begin{example}[Which clause of \texttt{LargeElim} applies]
  \label{def:iotashape-largeelim}
  \lean{Lean4Lean.Tests.IotaShape.resultSort, Lean4Lean.Tests.IotaShape.largeElimClause}
  \leanok
  \uses{def:ind-large-elim, def:ind-field-in-indices, def:meta-ofexpr}
  \contrib{iota}\srcloc{Lean4Lean/Tests/IotaShape.lean}{196}
  \thesisref{axioms.tex \S2.6.2, large elimination: the two reasons a type can eliminate largely.}

  \lead{In short} Decides which case of large elimination a block falls into.
  \begin{itemize}
  \item $\mathrm{resultSort}$: the sort a $\Pi$-telescope ends in.
  \item $\mathrm{largeElimClause}\;I$: $0$ if the result sort is never zero, $1$ single former with
        no constructor, $2$ single former with single constructor, $\mathrm{none}$ otherwise.
  \item Case $2$ also returns the fields the syntactic half leaves open
        (\texttt{VExpr.FieldInIndices}).
  \item \lead{Caveat} Whether those open fields are propositions is a typing condition,
        deliberately not decided here.
  \end{itemize}
\end{example}

\begin{example}[The block-level clauses, and which a nested block fails]
  \label{def:iotashape-blockfailures}
  \lean{Lean4Lean.Tests.IotaShape.recsOf, Lean4Lean.Tests.IotaShape.blockFailures, Lean4Lean.Tests.IotaShape.checkBlock, Lean4Lean.Tests.IotaShape.checkBlockRejected}
  \leanok
  \uses{def:ind-ctor-result, def:ind-ctor-positive, def:ind-rec-shape, def:ind-major-former, def:iotashape-largeelim, structure:ind-wf}
  \contrib{iota}\srcloc{Lean4Lean/Tests/IotaShape.lean}{224}
  \thesisref{axioms.tex \S2.6.1, inductive specifications; axioms.tex \S2.6.2, large elimination.}

  \lead{In short} Decides block-level clauses and names the ones a block fails.
  \begin{itemize}
  \item $\mathrm{recsOf}\;I$: $I.\mathrm{rec}$ plus any auxiliary \texttt{I.rec\_1},
        \texttt{I.rec\_2}, \dots
  \item $\mathrm{blockFailures}\;I$: names of failing clauses among
        \texttt{ctors\_params}/\texttt{result}/\texttt{positive}, \texttt{rec\_shape},
        \texttt{rec\_counts}, \texttt{recs\_over\_block}, \texttt{rules\_total}, syntactic
        \texttt{rules\_ctor}, and the large-elimination syntactic half.
  \item \texttt{checkBlock} demands an empty list; \texttt{checkBlockRejected} a non-empty one.
  \item \lead{Caveat} \texttt{checkBlockRejected} only asks that \emph{some} clause fail; which one
        is pinned only for \texttt{Tree} (\srcloc{Lean4Lean/Tests/IotaShape.lean}{584}), not the
        other eight nested fixtures.
  \end{itemize}
\end{example}

\begin{example}[The full battery per inductive type]
  \label{def:iotashape-checkall}
  \lean{Lean4Lean.Tests.IotaShape.checkAll, Lean4Lean.Tests.IotaShape.checkNested}
  \leanok
  \uses{def:iotashape-checkshapes, def:iotashape-checkmore, def:iotashape-checkiota, def:iotashape-blockfailures}
  \contrib{iota}\srcloc{Lean4Lean/Tests/IotaShape.lean}{313}
  \thesisref{no thesis counterpart; it packages the clauses of \ref{structure:ind-wf}.}

  \lead{In short} Runs the whole battery on one inductive type.
  \begin{itemize}
  \item $\mathrm{checkAll}\;I$: \texttt{checkTypesHaveRec}, \texttt{checkBlock}, then
        \texttt{checkShapes}, \texttt{checkMore}, \texttt{checkIotaAuto} on every recursor and
        rule.
  \item $\mathrm{checkNested}\;I$: demands the block be rejected, still comparing every iota rule
        with the kernel's.
  \end{itemize}
\end{example}

\begin{example}[Further shapes: nested, reflexive, indexed, small-eliminating]
  \label{def:iotashape-fixtures-nested}
  \lean{Lean4Lean.Tests.IotaShape.TreeP, Lean4Lean.Tests.IotaShape.T2, Lean4Lean.Tests.IotaShape.T3, Lean4Lean.Tests.IotaShape.MA, Lean4Lean.Tests.IotaShape.MB, Lean4Lean.Tests.IotaShape.T4, Lean4Lean.Tests.IotaShape.Refl, Lean4Lean.Tests.IotaShape.Le, Lean4Lean.Tests.IotaShape.Dep, Lean4Lean.Tests.IotaShape.TreeQ, Lean4Lean.Tests.IotaShape.Fn, Lean4Lean.Tests.IotaShape.EvI, Lean4Lean.Tests.IotaShape.OdI, Lean4Lean.Tests.IotaShape.PropLarge, Lean4Lean.Tests.IotaShape.Wrap, Lean4Lean.Tests.IotaShape.SigmaLike}
  \leanok
  \contrib{iota}\srcloc{Lean4Lean/Tests/IotaShape.lean}{339}
  \thesisref{axioms.tex \S2.6.1, inductive specifications.}

  \lead{In short} Sixteen further fixtures, one per remaining clause of \ref{structure:ind-wf}.
  \begin{itemize}
  \item Nesting through type formers or an indexed family, mutual-and-nested at once, reflexive
        constructors, an indexed \texttt{Prop}, a dependent structure, a mutual indexed
        \texttt{Type} block.
  \item Elimination boundary: two fixtures eliminate only into \texttt{Prop}; \texttt{Wrap}
        (single proof field) is subsingleton-eliminating and carries the extra universe parameter.
  \item \lead{Caveat} The docstring (\srcloc{Lean4Lean/Tests/IotaShape.lean}{336}) calls all three
        ``small-eliminating''; \texttt{Wrap} is in fact the large-eliminating case, only the
        wording is wrong.
  \end{itemize}
\end{example}

\section{Negative controls and the test battery}

Two declarations the elaborator refuses are written directly as \texttt{VExpr}/\texttt{VInductDecl}
literals, testing that the clauses reject what the kernel would reject.

\begin{example}[Strict positivity on hand-written constructor types]
  \label{thm:iotashape-positivity}
  \lean{Lean4Lean.Tests.IotaShape.badCtorType, Lean4Lean.Tests.IotaShape.goodCtorType, Lean4Lean.Tests.IotaShape.badCtorType_not_positive, Lean4Lean.Tests.IotaShape.goodCtorType_positive}
  \leanok
  \uses{def:ind-ctor-positive, fam:shapedecide-shapes}
  \contrib{iota}\srcloc{Lean4Lean/Tests/IotaShape.lean}{397}
  \thesisref{axioms.tex \S2.6.1, inductive specifications: the two kinds of constructor argument.}

  \lead{In short} A hand-written non-positive constructor type fails \texttt{CtorPositive}; a
  positive one passes.
  \begin{itemize}
  \item Over a fictitious former $\mathrm{Bad}$: $(\mathrm{Bad}\to\mathrm{False})\to\mathrm{Bad}$
        is \emph{not} $\mathtt{CtorPositive}\;[\mathrm{Bad}]\;0$.
  \item $(\mathrm{Nat}\to\mathrm{Bad})\to\mathrm{Bad}$ is. Both closed by \texttt{decide}, a
        control no kernel-declared type could otherwise exercise.
  \end{itemize}
\end{example}

\begin{example}[A two-constructor \texttt{Prop} asking for large elimination]
  \label{thm:iotashape-twoctorprop}
  \lean{Lean4Lean.Tests.IotaShape.TwoCtorProp.PT, Lean4Lean.Tests.IotaShape.TwoCtorProp.RT, Lean4Lean.Tests.IotaShape.TwoCtorProp.RV, Lean4Lean.Tests.IotaShape.TwoCtorProp.rhs1, Lean4Lean.Tests.IotaShape.TwoCtorProp.rhs2, Lean4Lean.Tests.IotaShape.TwoCtorProp.declB, Lean4Lean.Tests.IotaShape.TwoCtorProp.declB_wants_large, Lean4Lean.Tests.IotaShape.TwoCtorProp.declB_not_largeElimShape}
  \leanok
  \uses{def:ind-large-elim-shape, struct:vinductdecl, struct:vinductivetype, struct:vrecursor, struct:vrecrule, fam:shapedecide-shapes}
  \contrib{iota}\srcloc{Lean4Lean/Tests/IotaShape.lean}{411}
  \thesisref{axioms.tex \S2.6.2, large elimination: a \texttt{Prop} with two constructors is not
  subsingleton-eliminating.}

  \lead{In short} A hand-built unsound declaration that large elimination must reject.
  \begin{itemize}
  \item A \texttt{Prop} $P$ with two nullary constructors $t_1,t_2$ and a recursor eliminating
        into $\mathrm{Sort}\,(\mathrm{param}\;0)$: proof irrelevance identifies $t_1,t_2$, so with
        both iota rules one major premise would reduce to two different minors, collapsing
        definitional equality.
  \item Two \texttt{decide}d facts: the declaration \emph{does} ask for large elimination, and
        \texttt{VInductDecl.LargeElimShape} (\ref{def:ind-large-elim-shape}) refuses it.
  \end{itemize}
\end{example}

\begin{example}[The iota and shape test battery]
  \label{test:iotashape-run}
  \leanok
  \uses{def:iotashape-checkall, def:iotashape-checkiota, def:iotashape-ruleshapeat, def:iotashape-largeelim, def:iotashape-fixtures-direct, def:iotashape-fixtures-nested, def:ind-rec-shape, def:ind-rule-shape, def:ind-minor-for}
  \contrib{iota}\srcloc{Lean4Lean/Tests/IotaShape.lean}{454}
  \thesisref{axioms.tex \S2.6.3, the recursor; axioms.tex \S2.6.4, the computation rule (iota
  reduction).}

  \lead{In short} One \texttt{run\_meta} block running the whole battery, each check paired with a
  negative control.
  \begin{itemize}
  \item \texttt{checkShapes}/\texttt{checkTypesHaveRec} on ten recursors, each with a matching
        shape control (wrong telescope splits, rule counts, \texttt{recs\_elim} levels).
  \item Twelve \texttt{checkIota} comparisons; six \texttt{largeElimClause} checks; the documented
        \texttt{Tree} control.
  \item \texttt{checkAll} on forty-eight type formers, \texttt{checkNested} on nine nested blocks,
        including the \texttt{trproj} control at \srcloc{Lean4Lean/Tests/IotaShape.lean}{538}.
  \item \lead{Caveat} One \texttt{run\_meta} block: a failure reports only the first mismatch,
        leaving later assertions, negative controls included, unreached.
  \end{itemize}
\end{example}

\section{The kernel side: validating the expansion}

\texttt{Tests/ProjShape.lean} confronts the expansion with the real Lean kernel: for each field it
builds \ref{def:proj-fn} and \ref{def:proj-motive-body} and runs three kernel checks.

\begin{example}[Test structures]
  \label{def:projshape-fixtures}
  \lean{Lean4Lean.Tests.ProjShape.V3, Lean4Lean.Tests.ProjShape.Refl}
  \leanok
  \contrib{trproj}\srcloc{Lean4Lean/Tests/ProjShape.lean}{33}
  \thesisref{no counterpart; test fixtures.}

  \lead{In short} Two structures probing the expansion.
  \begin{itemize}
  \item \texttt{V3}: fields $n:\mathrm{Nat}$, $v:\mathrm{Fin}\;n$, $h:v.\mathrm{val}<n$ -- the last
        field's type mentions two earlier fields.
  \item \texttt{Refl}: single field $\mathrm{next}:\mathrm{Nat}\to\mathrm{Refl}$, a negative
        control whose minor premise carries an inductive-hypothesis binder, outside
        \ref{struct:trprojctor}'s scope.
  \end{itemize}
\end{example}

\begin{example}[Back-translation to kernel expressions]
  \label{fam:projshape-back-translation}
  \lean{Lean4Lean.Tests.ProjShape.toLevel, Lean4Lean.Tests.ProjShape.toExpr, Lean4Lean.Tests.ProjShape.binders}
  \leanok
  \uses{ind:vexpr, ind:vlevel, def:meta-ofexpr}
  \contrib{trproj}\srcloc{Lean4Lean/Tests/ProjShape.lean}{44}
  \thesisref{no counterpart; test scaffolding.}

  \lead{In short} Converts model terms back to kernel expressions for real-kernel checks.
  \begin{itemize}
  \item $\mathrm{toLevel}$: inverts \texttt{VLevel.ofLevel}. $\mathrm{toExpr}$: maps a
        \texttt{VExpr} to an \texttt{Expr} with dummy binder names. $\mathrm{binders}$: reads a
        type's leading $n$ binders.
  \item \lead{Caveat} $\mathrm{toLevel}$'s \texttt{ls.getD i .anonymous} silently defaults an
        out-of-range level index to the anonymous universe rather than erroring.
  \end{itemize}
\end{example}

\begin{example}[Kernel validation of the projection expansion]
  \label{def:projshape-checkproj}
  \lean{Lean4Lean.Tests.ProjShape.checkProj}
  \leanok
  \uses{def:proj-fn, def:proj-motive-body, def:proj-inst-pis, def:proj-binder-arity, def:pi-telescope, def:meta-ofexpr, def:kernel-infer-proj, fam:projshape-back-translation}
  \contrib{trproj}\srcloc{Lean4Lean/Tests/ProjShape.lean}{68}
  \thesisref{axioms.tex \S2.6.3, the recursor; typesys.tex \S3.1, the projection
  $\mathrm{inv}_x$ of $\mathrm{acc}\,x$ through $\mathrm{rec}_{\mathrm{acc}}$.}

  \lead{In short} Three independent checks that the real kernel accepts the expansion.
  \begin{itemize}
  \item The kernel accepts $\lambda\,ps\,(x{:}S\,ps).\;P_i\,x$ at type $\forall
        ps\,(x{:}S\,ps),\,\mathrm{projMotiveBody}$ as a definition.
  \item That type is \texttt{isDefEq} to the kernel's own $\mathrm{inferType}\,
        (\mathtt{.proj}\;S\;i\;x)$.
  \item $P_i\,(\mathrm{mk}\;ps\;fs)$ \texttt{whnf}-reduces \emph{syntactically} to $fs[i]$.
  \item \lead{Caveat} The second check is \texttt{isDefEq}, i.e.\ up to \emph{kernel} conversion
        (structure eta, projection reduction), which the model's \texttt{IsDefEq} deliberately
        lacks. $uss$ is built with \texttt{(VLevel.ofLevel lps (lvls j)).getD .zero}
        (\srcloc{Lean4Lean/Tests/ProjShape.lean}{78}), silently defaulting a mistyped level to
        zero.
  \end{itemize}
\end{example}

\begin{example}[The kernel-side test battery for projections]
  \label{test:projshape-run}
  \leanok
  \uses{def:projshape-checkproj, def:proj-binder-arity, def:projshape-fixtures}
  \contrib{trproj}\srcloc{Lean4Lean/Tests/ProjShape.lean}{121}
  \thesisref{axioms.tex \S2.6.3, the recursor.}

  \lead{In short} Runs \ref{def:projshape-checkproj} on eleven structure-and-field pairs.
  \begin{itemize}
  \item \texttt{Prod} 0,1; \texttt{Sigma} 0,1; \texttt{PSigma} 1; \texttt{Subtype} 0,1;
        \texttt{Fin} 1; \texttt{And} 1; \texttt{V3} 1,2 -- a constant motive, dependent motives at
        two universes, small elimination, and a two-projection chain.
  \item Two negative controls: a shared level list is rejected at \texttt{Sigma.snd};
        $\mathrm{binderArity?}\;1$ of \texttt{Refl.rec}'s minor is $2$, not $1$.
  \item \lead{Caveat} The \texttt{Sigma.snd} control (\srcloc{Lean4Lean/Tests/ProjShape.lean}{136})
        wraps the call in \texttt{try...catch \_ => pure false}, so any internal failure, not only
        a level mismatch, makes it pass.
  \end{itemize}
\end{example}

\begin{example}[The de Bruijn form of the \texttt{Sigma.snd} expansion]
  \label{test:projshape-sigmasnd}
  \lean{Lean4Lean.Tests.ProjShape.usS, Lean4Lean.Tests.ProjShape.uss, Lean4Lean.Tests.ProjShape.ps, Lean4Lean.Tests.ProjShape.Fs, Lean4Lean.Tests.ProjShape.structTy, Lean4Lean.Tests.ProjShape.P₀}
  \leanok
  \uses{def:proj-fn, def:proj-motive-body, def:proj-ty, def:proj-inst-pis}
  \contrib{trproj}\srcloc{Lean4Lean/Tests/ProjShape.lean}{150}
  \thesisref{Wtypes.tex \S5.1, the primitive $\Sigma$ rules $\pi_1\,p:\alpha$ and
  $\pi_2\,p:\beta[\pi_1\,p/x]$.}

  \lead{In short} Hand-checks the \texttt{Sigma.snd} expansion symbolically.
  \begin{itemize}
  \item Over $A:\mathrm{Type}\;u,\,B:A\to\mathrm{Type}\;v$: $P_0 =
        \mathrm{Sigma.rec}.\{u+1,u,v\}\;A\;B\;(\lambda \_.\,A)\;(\lambda a\,b.\,a)$.
  \item Five checks by \texttt{decide}/\texttt{simp}: \texttt{instPis} fails on a non-$\Pi$ head;
        $\mathrm{projFn}\dots 0=P_0$; the motive body of field $1$ over $\Gamma,x$ is $B\,(P_0\,x)$;
        $\mathrm{projFn}\dots 1$ uses the level list $\{v+1,u,v\}$; $\mathrm{projTy}\dots 1\;e=
        B\,(P_0\,e)$.
  \end{itemize}
\end{example}

\section{The model side: a structure with a constant motive}

\texttt{Tests/ProjInhabit.lean} builds inductive blocks by hand via \texttt{VEnv.addInduct}
(\ref{def:ind-add-induct}) and derives each projection's typing from \texttt{IsDefEq}
(\ref{def:isdefeq}) alone. \texttt{Plain} is the easy half: both fields share a type, so both
motives are constant.

\begin{example}[A hand-built two-field structure block]
  \label{def:projinhabit-plain-block}
  \lean{Lean4Lean.Tests.ProjInhabit.Plain.l1, Lean4Lean.Tests.ProjInhabit.Plain.Ac, Lean4Lean.Tests.ProjInhabit.Plain.Sc, Lean4Lean.Tests.ProjInhabit.Plain.mkc, Lean4Lean.Tests.ProjInhabit.Plain.ctorTy, Lean4Lean.Tests.ProjInhabit.Plain.motiveTy, Lean4Lean.Tests.ProjInhabit.Plain.mkSpine, Lean4Lean.Tests.ProjInhabit.Plain.minorTy, Lean4Lean.Tests.ProjInhabit.Plain.recTy, Lean4Lean.Tests.ProjInhabit.Plain.ruleRhs, Lean4Lean.Tests.ProjInhabit.Plain.decl, Lean4Lean.Tests.ProjInhabit.Plain.env0?, Lean4Lean.Tests.ProjInhabit.Plain.env0_isSome, Lean4Lean.Tests.ProjInhabit.Plain.env}
  \leanok
  \uses{def:ind-add-induct, struct:vinductdecl, struct:venv, def:venv-basic}
  \contrib{trproj}\srcloc{Lean4Lean/Tests/ProjInhabit.lean}{27}
  \thesisref{axioms.tex \S2.6.1, inductive specifications; axioms.tex \S2.6.3, the recursor.}

  \lead{In short} One former $S$, constructor $S.\mathrm{mk}:A\to A\to S$, one recursor with
  template $\lambda C\,m\,a\,b.\,m\,a\,b$.
  \begin{itemize}
  \item No universe variables, no parameters. Environment
        $\mathtt{VEnv.empty.addConst}\;A \mathbin{>\!\!>\!\!=} \mathtt{addInduct}\;\mathit{decl}$,
        shown to succeed by \texttt{rfl}.
  \item \lead{Caveat} \texttt{VEnv.addInduct} (\srcloc{Lean4Lean/Theory/Inductive.lean}{316})
        checks only name clashes and reduct closedness, never well-formedness: neither
        $\mathtt{VInductDecl.WF}$ nor $\mathit{env}.\mathtt{WF}$ is proved here.
  \end{itemize}
\end{example}

\begin{example}[Constant lookups and the registered iota rule]
  \label{fam:projinhabit-plain-lookups}
  \lean{Lean4Lean.Tests.ProjInhabit.Plain.hA, Lean4Lean.Tests.ProjInhabit.Plain.hS, Lean4Lean.Tests.ProjInhabit.Plain.hmk, Lean4Lean.Tests.ProjInhabit.Plain.hrec, Lean4Lean.Tests.ProjInhabit.Plain.hclosed, Lean4Lean.Tests.ProjInhabit.Plain.patKey, Lean4Lean.Tests.ProjInhabit.Plain.rhsPair, Lean4Lean.Tests.ProjInhabit.Plain.hpats}
  \leanok
  \uses{def:projinhabit-plain-block, def:pattern-iota-rhs, inductive:simple-pattern, struct:venv}
  \contrib{trproj}\srcloc{Lean4Lean/Tests/ProjInhabit.lean}{55}
  \thesisref{axioms.tex \S2.6.4: the computation rule.}

  \lead{In short} Basic \texttt{rfl} facts about the hand-built environment.
  \begin{itemize}
  \item Four constant lookups ($A$, $S$, $S.\mathrm{mk}$, $S.\mathrm{rec}$); closedness of the
        rule template by \texttt{decide}.
  \item $\mathit{env}.\mathrm{pats}\;\mathit{patKey}\;\mathit{rhsPair}$ at
        $\mathtt{SimplePattern.iota}\;S.\mathrm{rec}\;(0{+}1{+}1{+}0)\;S.\mathrm{mk}\;(0{+}2)$,
        exactly what \texttt{TrProjCtor.pat} demands.
  \end{itemize}
\end{example}

\begin{example}[The expansion terms for both fields]
  \label{def:projinhabit-plain-expansion}
  \lean{Lean4Lean.Tests.ProjInhabit.Plain.uss, Lean4Lean.Tests.ProjInhabit.Plain.recC, Lean4Lean.Tests.ProjInhabit.Plain.M0, Lean4Lean.Tests.ProjInhabit.Plain.sel0, Lean4Lean.Tests.ProjInhabit.Plain.sel1, Lean4Lean.Tests.ProjInhabit.Plain.minorTyM}
  \leanok
  \uses{def:proj-fn, def:proj-motive-body, def:projinhabit-plain-block}
  \contrib{trproj}\srcloc{Lean4Lean/Tests/ProjInhabit.lean}{70}
  \thesisref{no counterpart; an instance of the expansion.}

  \lead{In short} The constant motive and two selectors, checked against \ref{def:proj-fn}.
  \begin{itemize}
  \item $M_0=\lambda x{:}S.\,A$, $\mathit{sel}_0=\lambda a\,b.\,a$, $\mathit{sel}_1=\lambda
        a\,b.\,b$.
  \item By \texttt{decide}: $\mathrm{projFn}\;S\;[]\;\mathit{uss}\;[]\;[A,A]\;i =
        S.\mathrm{rec}\,[1]\;M_0\;\mathit{sel}_i$ and $\mathrm{projMotiveBody}\dots i=A$, for
        $i=0,1$.
  \end{itemize}
\end{example}

\begin{example}[Typing derivation of the constant-motive expansion]
  \label{fam:projinhabit-plain-typing}
  \lean{Lean4Lean.Tests.ProjInhabit.Plain.lvl1WF, Lean4Lean.Tests.ProjInhabit.Plain.hAt, Lean4Lean.Tests.ProjInhabit.Plain.hSt, Lean4Lean.Tests.ProjInhabit.Plain.hmkt, Lean4Lean.Tests.ProjInhabit.Plain.hrect, Lean4Lean.Tests.ProjInhabit.Plain.hM, Lean4Lean.Tests.ProjInhabit.Plain.hspine, Lean4Lean.Tests.ProjInhabit.Plain.hbody0, Lean4Lean.Tests.ProjInhabit.Plain.hselnat0, Lean4Lean.Tests.ProjInhabit.Plain.hselnat1, Lean4Lean.Tests.ProjInhabit.Plain.hminorEq, Lean4Lean.Tests.ProjInhabit.Plain.hsel0, Lean4Lean.Tests.ProjInhabit.Plain.hsel1, Lean4Lean.Tests.ProjInhabit.Plain.hstep1, Lean4Lean.Tests.ProjInhabit.Plain.hbetax, Lean4Lean.Tests.ProjInhabit.Plain.htyEq, Lean4Lean.Tests.ProjInhabit.Plain.hProjFn0, Lean4Lean.Tests.ProjInhabit.Plain.hProjFn1}
  \leanok
  \uses{def:isdefeq, def:projinhabit-plain-expansion, def:projinhabit-plain-block}
  \contrib{trproj}\srcloc{Lean4Lean/Tests/ProjInhabit.lean}{84}
  \thesisref{steps (a) and (c) of the derivation sketched at
  \srcloc{Lean4Lean/Theory/Proj.lean}{38}; axioms.tex \S2.6.3, the recursor.}

  \lead{In short} Seventeen term-mode \texttt{IsDefEq} derivations culminating in $P_i$'s typing.
  \begin{itemize}
  \item Ends at $\mathit{env}\vdash_\Gamma S.\mathrm{rec}\,[1]\;M_0\;\mathit{sel}_i:\forall
        x{:}S,\,A$.
  \item \lead{Idea} Type the motive (\texttt{lamDF}); type the selector at its natural type;
        convert to the minor type by $\beta$ under the field binders; apply (\texttt{appDF});
        convert the result by $\beta$ at the major (\texttt{defeqDF}).
  \end{itemize}
\end{example}

\begin{example}[\texttt{TrProjCtor} is inhabited for field 0]
  \label{thm:projinhabit-plain-inhab0}
  \lean{Lean4Lean.Tests.ProjInhabit.Plain.inhab0}
  \leanok
  \uses{struct:trprojctor, fam:projinhabit-plain-typing, fam:projinhabit-plain-lookups, def:proj-binder-arity, def:proj-inst-pis}
  \contrib{trproj}\srcloc{Lean4Lean/Tests/ProjInhabit.lean}{151}
  \thesisref{typesys.tex \S3.1, the projection $\mathrm{inv}_x$ of $\mathrm{acc}\,x$ through
  $\mathrm{rec}_{\mathrm{acc}}$.}

  \lead{In short} \hl{\texttt{TrProjCtor} is not vacuous}: every clause is discharged for field
  $0$.
  \begin{itemize}
  \item $\mathtt{TrProjCtor}\;\mathit{env}\;0\;[S]\;S\;0\;(\mathrm{bvar}\;0)\;
        (S.\mathrm{rec}\,[1]\;M_0\;\mathit{sel}_0\;(\mathrm{bvar}\;0))\;S.\mathrm{mk}\;[]\;
        \mathit{uss}\;[]\;0\;[A,A]$ holds.
  \item An anonymous constructor: \texttt{pat} from \ref{fam:projinhabit-plain-lookups};
        \texttt{ctor}, \texttt{minor\_arity}, \texttt{eq} by \texttt{rfl}; \texttt{field\_lt} by
        \texttt{decide}; \texttt{major\_ty} by \texttt{bvar}; \texttt{fn\_ty} from
        \ref{fam:projinhabit-plain-typing}.
  \item \lead{Caveat} Shows \texttt{TrProjCtor} satisfiable, not that a kernel-accepted projection
        yields one -- that bridge, \texttt{inferProj.WF\_struct}, is \texttt{sorry}.
  \end{itemize}
\end{example}

\begin{example}[The same for field 1, and the packaged \texttt{TrProj}]
  \label{thm:projinhabit-plain-inhab1}
  \lean{Lean4Lean.Tests.ProjInhabit.Plain.inhab1, Lean4Lean.Tests.ProjInhabit.Plain.trProj0}
  \leanok
  \uses{thm:projinhabit-plain-inhab0, def:trproj}
  \contrib{trproj}\srcloc{Lean4Lean/Tests/ProjInhabit.lean}{164}
  \thesisref{no counterpart; a second instance of the same shape.}

  \lead{In short} The same for field $1$, packaged as \texttt{TrProj}.
  \begin{itemize}
  \item \texttt{inhab1}: structure update of \texttt{inhab0} at \texttt{field\_lt},
        \texttt{fn\_ty}, \texttt{eq} (field $1$'s type does not mention field $0$, so its motive is
        also constant).
  \item \texttt{trProj0} packs \texttt{inhab0} into the existential \texttt{TrProj}
        (\ref{def:trproj}).
  \end{itemize}
\end{example}

\section{The model side: a structure with a dependent motive}

\texttt{Dependent} types the second selector at $\forall a\,b,\,M_{21}\,(\mathrm{mk}\;a\;b)$, which
needs $B\,a\equiv B\,(P_0\,(\mathrm{mk}\;a\;b))$: the first projection computing on the generic
spine, where the registered iota rule actually fires.

\begin{example}[A hand-built Sigma-shaped block]
  \label{def:projinhabit-dep-block}
  \lean{Lean4Lean.Tests.ProjInhabit.Dependent.Bc, Lean4Lean.Tests.ProjInhabit.Dependent.Bx, Lean4Lean.Tests.ProjInhabit.Dependent.BTy, Lean4Lean.Tests.ProjInhabit.Dependent.ctorTy2, Lean4Lean.Tests.ProjInhabit.Dependent.motiveTy2, Lean4Lean.Tests.ProjInhabit.Dependent.motiveTy2', Lean4Lean.Tests.ProjInhabit.Dependent.minorTy2, Lean4Lean.Tests.ProjInhabit.Dependent.recTy2, Lean4Lean.Tests.ProjInhabit.Dependent.ruleRhs2, Lean4Lean.Tests.ProjInhabit.Dependent.decl2, Lean4Lean.Tests.ProjInhabit.Dependent.env0?, Lean4Lean.Tests.ProjInhabit.Dependent.env, Lean4Lean.Tests.ProjInhabit.Dependent.env0_isSome}
  \leanok
  \uses{def:ind-add-induct, struct:vinductdecl, def:venv-basic}
  \contrib{trproj}\srcloc{Lean4Lean/Tests/ProjInhabit.lean}{181}
  \thesisref{axioms.tex \S2.6.1, inductive specifications.}

  \lead{In short} The dependent counterpart: $S_2.\mathrm{mk}:(a{:}A)\to B\,a\to S_2$.
  \begin{itemize}
  \item Over an environment also declaring $B:A\to\mathrm{Sort}\;1$; the second field's type
        mentions the first, forcing the dependent motive $\lambda x.\,B\,(P_0\,x)$.
  \end{itemize}
\end{example}

\begin{example}[Constant lookups for the dependent block]
  \label{fam:projinhabit-dep-lookups}
  \lean{Lean4Lean.Tests.ProjInhabit.Dependent.hA, Lean4Lean.Tests.ProjInhabit.Dependent.hB, Lean4Lean.Tests.ProjInhabit.Dependent.hS2, Lean4Lean.Tests.ProjInhabit.Dependent.hmk2, Lean4Lean.Tests.ProjInhabit.Dependent.hrec2, Lean4Lean.Tests.ProjInhabit.Dependent.hclosed2, Lean4Lean.Tests.ProjInhabit.Dependent.patKey2, Lean4Lean.Tests.ProjInhabit.Dependent.rhsPair2, Lean4Lean.Tests.ProjInhabit.Dependent.hpats2}
  \leanok
  \uses{def:projinhabit-dep-block, def:pattern-iota-rhs, inductive:simple-pattern, struct:venv}
  \contrib{trproj}\srcloc{Lean4Lean/Tests/ProjInhabit.lean}{215}
  \thesisref{axioms.tex \S2.6.4: the computation rule.}

  \lead{In short} As \ref{fam:projinhabit-plain-lookups}, for the dependent block.
  \begin{itemize}
  \item Five \texttt{rfl} lookups; closedness by \texttt{decide}; registered entry at
        $\mathtt{iota}\;S_2.\mathrm{rec}\;2\;S_2.\mathrm{mk}\;2$.
  \end{itemize}
\end{example}

\begin{example}[The dependent expansion terms]
  \label{def:projinhabit-dep-expansion}
  \lean{Lean4Lean.Tests.ProjInhabit.Dependent.uss, Lean4Lean.Tests.ProjInhabit.Dependent.rec2C, Lean4Lean.Tests.ProjInhabit.Dependent.M20, Lean4Lean.Tests.ProjInhabit.Dependent.sel20, Lean4Lean.Tests.ProjInhabit.Dependent.sel21, Lean4Lean.Tests.ProjInhabit.Dependent.P0', Lean4Lean.Tests.ProjInhabit.Dependent.M21, Lean4Lean.Tests.ProjInhabit.Dependent.P1', Lean4Lean.Tests.ProjInhabit.Dependent.minorTyM20, Lean4Lean.Tests.ProjInhabit.Dependent.minorTyM21}
  \leanok
  \uses{def:proj-fn, def:proj-motive-body, def:proj-motive-body-of, def:projinhabit-dep-block}
  \contrib{trproj}\srcloc{Lean4Lean/Tests/ProjInhabit.lean}{231}
  \thesisref{Wtypes.tex \S5.1, $\pi_2\,p:\beta[\pi_1\,p/x]$.}

  \lead{In short} $P_0$ constant, $P_1$ \hl{genuinely dependent}.
  \begin{itemize}
  \item $M_{20}=\lambda x{:}S_2.\,A$; $M_{21}=\lambda x.\,B\,(P_0\,x)$.
  \item By \texttt{decide}: \ref{def:proj-fn} and \ref{def:proj-motive-body} at $i=0,1$ equal
        these terms -- the de Bruijn bookkeeping of \ref{def:proj-motive-body-of} really produces
        $\lambda x.\,B\,(P_0\,x)$.
  \end{itemize}
\end{example}

\begin{example}[The level-instantiated iota template]
  \label{def:projinhabit-dep-template}
  \lean{Lean4Lean.Tests.ProjInhabit.Dependent.T, Lean4Lean.Tests.ProjInhabit.Dependent.bodyT, Lean4Lean.Tests.ProjInhabit.Dependent.inner3, Lean4Lean.Tests.ProjInhabit.Dependent.T1, Lean4Lean.Tests.ProjInhabit.Dependent.T2, Lean4Lean.Tests.ProjInhabit.Dependent.T3, Lean4Lean.Tests.ProjInhabit.Dependent.T5}
  \leanok
  \uses{def:instl, def:projinhabit-dep-block}
  \contrib{trproj}\srcloc{Lean4Lean/Tests/ProjInhabit.lean}{248}
  \thesisref{axioms.tex \S2.6.4: the computation rule.}

  \lead{In short} Names the intermediate terms of a six-step $\beta$-reduction.
  \begin{itemize}
  \item $T$: the rule template level-instantiated at $[l_1]$ (\texttt{decide}).
  \item $\mathit{bodyT}$, $\mathit{inner3}$, $T_1,T_2,T_3,T_5$: intermediate terms of $T\;
        M_{20}\;\mathit{sel}_{20}\;a\;b$'s reduction to $a$.
  \end{itemize}
\end{example}

\begin{example}[Basic typings for the dependent block]
  \label{fam:projinhabit-dep-basic-typing}
  \lean{Lean4Lean.Tests.ProjInhabit.Dependent.lvl1WF, Lean4Lean.Tests.ProjInhabit.Dependent.hAt, Lean4Lean.Tests.ProjInhabit.Dependent.hBt, Lean4Lean.Tests.ProjInhabit.Dependent.hS2t, Lean4Lean.Tests.ProjInhabit.Dependent.hmk2t, Lean4Lean.Tests.ProjInhabit.Dependent.hrec2t, Lean4Lean.Tests.ProjInhabit.Dependent.hBxt, Lean4Lean.Tests.ProjInhabit.Dependent.hmk2spine, Lean4Lean.Tests.ProjInhabit.Dependent.hM20, Lean4Lean.Tests.ProjInhabit.Dependent.hbody20, Lean4Lean.Tests.ProjInhabit.Dependent.hsel20nat, Lean4Lean.Tests.ProjInhabit.Dependent.hminorEq20, Lean4Lean.Tests.ProjInhabit.Dependent.hsel20M, Lean4Lean.Tests.ProjInhabit.Dependent.hstep20, Lean4Lean.Tests.ProjInhabit.Dependent.hbetax20, Lean4Lean.Tests.ProjInhabit.Dependent.hM21body, Lean4Lean.Tests.ProjInhabit.Dependent.hM21, Lean4Lean.Tests.ProjInhabit.Dependent.hredex}
  \leanok
  \uses{def:isdefeq, def:projinhabit-dep-expansion}
  \contrib{trproj}\srcloc{Lean4Lean/Tests/ProjInhabit.lean}{262}
  \thesisref{no counterpart; instance-level typing bookkeeping.}

  \lead{In short} Constant typings and the generic constructor spine's typing.
  \begin{itemize}
  \item $\Gamma,a{:}A,b{:}B\,a \vdash \mathrm{mk}\;a\;b : S_2$; the pieces of $P_0$'s typing; the
        typing of $M_{21}$; well-typedness of the redex $P_0\,(\mathrm{mk}\;a\;b)$.
  \item Term-mode \texttt{IsDefEq} constructors, as in \texttt{Plain}.
  \end{itemize}
\end{example}

\begin{example}[Typing of the first projection function]
  \label{thm:projinhabit-dep-hp0ty}
  \lean{Lean4Lean.Tests.ProjInhabit.Dependent.hP0ty}
  \leanok
  \uses{fam:projinhabit-dep-basic-typing, def:isdefeq}
  \contrib{trproj}\srcloc{Lean4Lean/Tests/ProjInhabit.lean}{316}
  \thesisref{step (c) of the derivation at \srcloc{Lean4Lean/Theory/Proj.lean}{49}.}

  \lead{In short} $P_0$'s typing, the prerequisite for field $1$'s dependent motive.
  \begin{itemize}
  \item $\mathit{env} \vdash_\Gamma P_0 : \forall x{:}S_2,\,A$, in any context.
  \item \lead{Idea} \texttt{defeqDF} of \texttt{forallEDF hS2t hbetax20} onto
        \texttt{appDF hstep20 hsel20M}.
  \end{itemize}
\end{example}

\begin{example}[The iota pattern matches the spine]
  \label{thm:projinhabit-dep-hmatch}
  \lean{Lean4Lean.Tests.ProjInhabit.Dependent.hmatch}
  \leanok
  \uses{inductive:pattern-matches, def:pattern-iota-rhs, fam:projinhabit-dep-lookups}
  \contrib{trproj}\srcloc{Lean4Lean/Tests/ProjInhabit.lean}{337}
  \thesisref{axioms.tex \S2.6.4: the computation rule.}

  \lead{In short} The iota pattern matches $P_0$ applied to the generic constructor spine.
  \begin{itemize}
  \item An assignment $m_2$ with $\mathit{patKey}_2.\mathrm{Matches}\;(P_0\,(\mathrm{mk}\;a\;b))\;
        [l_1]\;m_2$, reduct $T\;M_{20}\;\mathit{sel}_{20}\;(\mathrm{bvar}\;1)\;(\mathrm{bvar}\;0)$.
  \item \lead{Idea} Explicit derivation \texttt{.app (.var (.var .const)) (.var (.var .const))},
        closed by \texttt{rfl}.
  \end{itemize}
\end{example}

\begin{example}[The iota step on the constructor spine]
  \label{thm:projinhabit-dep-hiota}
  \lean{Lean4Lean.Tests.ProjInhabit.Dependent.hiota}
  \leanok
  \uses{rule:isdefeq-pat, thm:projinhabit-dep-hmatch, fam:projinhabit-dep-basic-typing, fam:projinhabit-dep-lookups}
  \contrib{trproj}\srcloc{Lean4Lean/Tests/ProjInhabit.lean}{344}
  \thesisref{axioms.tex \S2.6.4: the computation rule.}

  \lead{In short} \hl{The model's iota rule actually fires} here.
  \begin{itemize}
  \item $\mathit{env}\vdash_{\Gamma,a{:}A,b{:}B\,a} P_0\,(\mathrm{mk}\;a\;b)\equiv T\;M_{20}\;
        \mathit{sel}_{20}\;a\;b:A$.
  \item \lead{Idea} \texttt{VEnv.IsDefEq.pat} (\ref{rule:isdefeq-pat}) at the match of
        \ref{thm:projinhabit-dep-hmatch}, the redex's typing, and empty side conditions,
        rewritten by the computed reduct.
  \end{itemize}
\end{example}

\begin{example}[Step-by-step reduction of the iota reduct]
  \label{fam:projinhabit-dep-beta-chain}
  \lean{Lean4Lean.Tests.ProjInhabit.Dependent.hbvarC, Lean4Lean.Tests.ProjInhabit.Dependent.hCapp, Lean4Lean.Tests.ProjInhabit.Dependent.hminorTy2wf, Lean4Lean.Tests.ProjInhabit.Dependent.hbvarM, Lean4Lean.Tests.ProjInhabit.Dependent.hmSpine1, Lean4Lean.Tests.ProjInhabit.Dependent.hmSpine2, Lean4Lean.Tests.ProjInhabit.Dependent.hbodyT, Lean4Lean.Tests.ProjInhabit.Dependent.hbeta1, Lean4Lean.Tests.ProjInhabit.Dependent.hc1, Lean4Lean.Tests.ProjInhabit.Dependent.hbvarM', Lean4Lean.Tests.ProjInhabit.Dependent.hmSpine1', Lean4Lean.Tests.ProjInhabit.Dependent.hmSpine2', Lean4Lean.Tests.ProjInhabit.Dependent.hinner3, Lean4Lean.Tests.ProjInhabit.Dependent.hbeta2, Lean4Lean.Tests.ProjInhabit.Dependent.hc2, Lean4Lean.Tests.ProjInhabit.Dependent.hinnerT2, Lean4Lean.Tests.ProjInhabit.Dependent.hbeta3, Lean4Lean.Tests.ProjInhabit.Dependent.hc3, Lean4Lean.Tests.ProjInhabit.Dependent.hbeta4, Lean4Lean.Tests.ProjInhabit.Dependent.hbeta5, Lean4Lean.Tests.ProjInhabit.Dependent.hc5, Lean4Lean.Tests.ProjInhabit.Dependent.hbeta6}
  \leanok
  \uses{def:isdefeq, def:projinhabit-dep-template}
  \contrib{trproj}\srcloc{Lean4Lean/Tests/ProjInhabit.lean}{353}
  \thesisref{axioms.tex \S2.6.4: the computation rule.}

  \lead{In short} Twenty-two lemmas typing every intermediate term of $T\;M_{20}\;
  \mathit{sel}_{20}\;a\;b\to^{*}a$.
  \begin{itemize}
  \item Each $\beta$ step types its body under its binder; each congruence re-applies
        \texttt{appDF} at the current type.
  \item Only \texttt{beta}, \texttt{appDF}, \texttt{lamDF}, \texttt{forallEDF}, \texttt{bvar}; no
        tactics.
  \item \lead{Caveat} \texttt{VEnv.IsDefEqU.betaN} (\srcloc{Lean4Lean/Verify/Environment/Lemmas.lean}{994})
        would shorten this chain but needs \texttt{VEnv.WF}, absent from these hand-built
        environments.
  \end{itemize}
\end{example}

\begin{example}[The first projection computes on the spine]
  \label{thm:projinhabit-dep-hproj0red}
  \lean{Lean4Lean.Tests.ProjInhabit.Dependent.hproj0red}
  \leanok
  \uses{thm:projinhabit-dep-hiota, fam:projinhabit-dep-beta-chain, def:isdefeq}
  \contrib{trproj}\srcloc{Lean4Lean/Tests/ProjInhabit.lean}{485}
  \thesisref{axioms.tex \S2.6.4: the computation rule.}

  \lead{In short} The full iota-plus-$\beta$ chain to the field's value.
  \begin{itemize}
  \item $\mathit{env}\vdash_{\Gamma,a{:}A,b{:}B\,a} P_0\,(\mathrm{mk}\;a\;b)\equiv a:A$.
  \item \lead{Idea} Seven \texttt{trans} steps chaining \ref{thm:projinhabit-dep-hiota}, the
        congruences \texttt{hc1,hc2,hc3,hc5}, and the $\beta$ steps \texttt{hbeta4,hbeta6}. The
        model-side counterpart of the kernel \texttt{whnf} check in \ref{def:projshape-checkproj}.
  \end{itemize}
\end{example}

\begin{example}[The typing premise for the dependent field]
  \label{thm:projinhabit-dep-hp1ty}
  \lean{Lean4Lean.Tests.ProjInhabit.Dependent.hP1ty, Lean4Lean.Tests.ProjInhabit.Dependent.hsel21nat, Lean4Lean.Tests.ProjInhabit.Dependent.hbody21, Lean4Lean.Tests.ProjInhabit.Dependent.hbodyEq21, Lean4Lean.Tests.ProjInhabit.Dependent.hminorEq21, Lean4Lean.Tests.ProjInhabit.Dependent.hsel21M, Lean4Lean.Tests.ProjInhabit.Dependent.hstep21, Lean4Lean.Tests.ProjInhabit.Dependent.hbetax21}
  \leanok
  \uses{thm:projinhabit-dep-hproj0red, def:isdefeq, def:projinhabit-dep-expansion}
  \contrib{trproj}\srcloc{Lean4Lean/Tests/ProjInhabit.lean}{531}
  \thesisref{step (b) of the derivation at \srcloc{Lean4Lean/Theory/Proj.lean}{45};
  Wtypes.tex \S5.1, $\pi_2\,p:\beta[\pi_1\,p/x]$.}

  \lead{In short} $P_1$ is typable using only $\beta$ and the registered iota rule.
  \begin{itemize}
  \item $\mathit{env}\vdash_\Gamma P_1 : \forall x{:}S_2,\;B\,(P_0\,x)$.
  \item \lead{Idea} Convert the selector's natural type $\forall a\,b,\;B\,a$ to the minor type via
        $B\,a\equiv B\,(P_0\,(\mathrm{mk}\;a\;b))$ (\ref{thm:projinhabit-dep-hproj0red}, under
        \texttt{appDF}); \texttt{forallEDF}/\texttt{defeqDF} give the minor typing; \texttt{appDF}
        applies the recursor; a final $\beta$ \texttt{defeqDF} gives the dependent typing -- no
        structure eta, no injectivity.
  \end{itemize}
\end{example}

\begin{example}[\texttt{TrProjCtor} and \texttt{TrProj} for the dependent block]
  \label{thm:projinhabit-dep-inhab}
  \lean{Lean4Lean.Tests.ProjInhabit.Dependent.inhabDep0, Lean4Lean.Tests.ProjInhabit.Dependent.inhabDep1, Lean4Lean.Tests.ProjInhabit.Dependent.trProjDep1}
  \leanok
  \uses{struct:trprojctor, def:trproj, thm:projinhabit-dep-hp0ty, thm:projinhabit-dep-hp1ty, fam:projinhabit-dep-lookups}
  \contrib{trproj}\srcloc{Lean4Lean/Tests/ProjInhabit.lean}{538}
  \thesisref{typesys.tex \S3.1, the projection $\mathrm{inv}_x$ of $\mathrm{acc}\,x$ through
  $\mathrm{rec}_{\mathrm{acc}}$.}

  \lead{In short} \texttt{TrProjCtor} inhabited for both fields of the dependent block.
  \begin{itemize}
  \item \texttt{inhabDep0}, \texttt{inhabDep1}:
        $\mathtt{TrProjCtor}\;\mathit{env}\;0\;[S_2]\;S_2\;i\;(\mathrm{bvar}\;0)\;
        (P_i\,(\mathrm{bvar}\;0))\;S_2.\mathrm{mk}\;[]\;\mathit{uss}\;[]\;0\;[A, B\,a]$, the second
        with the dependent motive.
  \item \texttt{trProjDep1} packs the latter into \texttt{TrProj}.
  \item \lead{Caveat} Together with \ref{thm:projinhabit-plain-inhab0}/\texttt{1}, the only
        evidence \texttt{TrProjCtor} is satisfiable; neither namespace has parameters, a
        non-trivial $usS$, per-field $uss$, or more than two fields.
  \end{itemize}
\end{example}

\begin{example}[Axiom profile of the witnesses]
  \label{test:projinhabit-axioms}
  \leanok
  \uses{thm:tr-env-proj-defeq, thm:wf-patsstrong, thm:projinhabit-dep-inhab, thm:projinhabit-plain-inhab0}
  \contrib{trproj}\srcloc{Lean4Lean/Tests/ProjInhabit.lean}{562}
  \thesisref{no counterpart; a trust-boundary guard.}

  \lead{In short} Pins the axiom profile of the four witnesses and of
  \ref{thm:tr-env-proj-defeq}.
  \begin{itemize}
  \item The four witnesses (\texttt{Plain.inhab0/1}, \texttt{Dependent.inhabDep0/1}) depend only
        on \texttt{[propext, Quot.sound]}: sorry-free.
  \item \texttt{Lean4Lean.TrEnv.proj\_defeq} (\ref{thm:tr-env-proj-defeq}) additionally depends on
        \texttt{sorryAx} plus three \texttt{PersistentHashMap} axioms: this, not a literal
        \texttt{sorry}, is where the module records the open obligation.
  \end{itemize}
  \axfoot{sorryAx, reached through \texttt{VEnv.WF.orderedStrong} from \texttt{VEnv.WF.patsStrong}
  (\srcloc{Lean4Lean/Theory/Typing/EnvLemmas.lean}{334}); also propext, Classical.choice,
  Quot.sound and three \texttt{Lean.PersistentHashMap} axioms.}
\end{example}

\section*{Caveats}
\label{sec:proj-review}

\begin{itemize}
\item The bridge from the kernel is open: \texttt{inferProj.WF\_struct} (Chapter~\ref{chap:trexpr})
is \texttt{sorry}, so no kernel-accepted projection is shown to yield a \texttt{TrProjCtor}.
\item Structure eta is outside what this representation can express.
\item \texttt{VExpr.projTy} (\ref{def:proj-ty}) is unused; the identity connecting it to
\ref{def:proj-motive-body} is asserted in its docstring but never proved.
\item \texttt{VEnv.addInduct} checks no well-formedness; the two hand-built witness environments
are never shown \texttt{WF}, so they establish satisfiability, not well-formed satisfiability.
\item The witnesses cover no parameters, no non-trivial $usS$, and a constant per-field level
list; \ref{fam:proj-fn-subst} meets no non-trivial instance in any test.
\item Thesis references (\texttt{inv\_x}, $\pi_2$) are borrowed idioms: their source systems are
not the ones this development formalises.
\end{itemize}
